\section{Whole polynomial projective comparison under coefficient isomorphisms}
\label{sec:polynomial-projective-coefficient-isomorphism}

\begin{theorem}[Whole coefficient isomorphism]
\label{thm:polynomial-projective-coefficient-isomorphism}
\lean{EnvelopeProperty.PolynomialProjectiveCoefficientIso.gradedMap, EnvelopeProperty.PolynomialProjectiveCoefficientIso.iso, EnvelopeProperty.PolynomialProjectiveCoefficientIso.iso_structure}
\leanok
\uses{thm:fc-symmetric-projective-embedding}
Any actual isomorphism $R\simeq S$ of commutative coefficient rings
induces an isomorphism of the entire original polynomial projective
schemes with the same variables. The whole original structural
morphisms commute with the spectrum of that same coefficient
isomorphism.
\end{theorem}

\begin{proof}
Apply the original coefficient isomorphism to every coefficient of
the original polynomial. The actual inverse coefficient map gives
its inverse on the entire polynomial ring. Both maps preserve every
homogeneous degree because they change no exponents. Surjectivity
gives the actual irrelevant-ideal inequalities. Native projective
functoriality and the original ring inverse identities then give
both whole scheme inverse identities. The degree-zero restriction
of the polynomial map sends the constant polynomial $r$ to the
constant polynomial given by the same coefficient isomorphism.
The resulting degree-zero square, followed by the original
projective-to-degree-zero structural map, proves the claimed
whole coefficient diagram.
\end{proof}
